\lecture{5}{September 8, 2025}{Differential Analysis I: Mass Conservation}

\section{Introduction to Differential Analysis of Fluid Motion}
In the previous \Cref{sec:intconvol} the basic equations in integral form for a control volume were introduced. These are useful when we are interested in the overall behaviour of a fluid field and its effects on its surroundings. To determine what happens at each point of the flow field we must ude the differential forms of the equations of motion.

\subsection{Conservation of Mass}

\subsubsection{Rectangular Coordinate System}
To find the differential equation for conservation of mass in a rectangular coordinate system we choose a control volume that is an infinitesimal cube with sides of length $\mathrm{d}x, \mathrm{d}y, \mathrm{d}z$ as shown on \Cref{fig:f5_1}. The density at the center, $O$, of the control volume is assumed to be $\rho$ and the velocity there is assumed to be $\textbf{V} = \hat{i} u + \hat{j} v + \hat{k} w$.

\begin{figure} [ht]
  \centering
  \includegraphics[width=0.6\linewidth]{./figures/f5_1.png}
  \caption{Differential control volume in rectangular coordinates.}
  \label{fig:f5_1}
\end{figure}

Now, to evaluate the properties at each of the six faces of the control series we will use a Taylor series expansion about point $O$. E.g. for the right face:
\[ 
\rho)_{x + \frac{\mathrm{d}x}{2}} = \rho + \left( \frac{\partial \rho}{\partial x} \right) \frac{\mathrm{d}x}{2} + \left( \frac{\partial^2 \rho}{\partial x^2} \right) \frac{1}{2!} \left( \frac{\mathrm{d}x}{2} \right)^2 + \ldots
.\]
By neglecting higher-order terms, we can write:
\[ 
\rho)_{x + \frac{\mathrm{d}x}{2}} = \rho + \left( \frac{\partial \rho}{\partial x} \right) \frac{\mathrm{d}x}{2}
.\]
Likewise for the velocity we can write:
\[ 
u )_{x + \frac{\mathrm{d}x}{2}} = u + \left( \frac{\partial u}{\partial x} \right) \frac{\mathrm{d}x}{2}
\]
where $\rho$, $u$, $\frac{\partial p}{\partial x}$, and $\frac{\partial u}{\partial x}$ are all evaluated at point $O$. Following the same procedure for the left hand side we get:
\begin{align*}
  \rho)_{x - \frac{\mathrm{d}x}{2}} &= \rho - \left( \frac{\partial p}{\partial x} \right) \frac{\mathrm{d}x}{2} \\
  u)_{x - \frac{\mathrm{d}x}{2}} &= u - \left( \frac{\partial u}{\partial x} \right) \frac{\mathrm{d}x}{2}
.\end{align*}
Similar expressions for $\rho$ and $v$ for the front and back faces and $\rho$ and $w$ for the top and bottom faces of the infinitesimal control volume can be made. These can then be used to evaluate the surface integral in \Cref{eq:convolconmas} (i.e. we can find the net flux of mass out of the control volume). \Cref{eq:convolconmas} states
\[ 
\frac{\partial }{\partial t} \int_{\mathrm{CV}} \rho \, \mathrm{d}V + \int_{\mathrm{CS}} \rho \textbf{V} \cdot \mathrm{d}\textbf{A} = 0
.\]
If we assume that the velocity components $u$, $v$, and $w$ are positive in the $x$, $y$, and $z$ directions, respectively, that the area normal is positive out of the cube and that higher order terms can be neglected as we find the limit of $\mathrm{d}x$, $\mathrm{d}y$, and $\mathrm{d}z$ as they go to zero, then we can solve the surface integral to be:
\[ 
\int_{\mathrm{CS}} \rho \textbf{V} \cdot \mathrm{d}\textbf{A} = \left( \frac{\partial \rho u}{\partial x} + \frac{\partial \rho v}{\partial y} + \frac{\partial \rho w}{\partial z} \right) \, \mathrm{d}x \, \mathrm{d}y \, \mathrm{d}z
.\]
To complete \Cref{eq:convolconmas} we need to evaluate the volume integral:
\[ 
\frac{\partial }{\partial t} \int_{\mathrm{CV}} \rho \, \mathrm{d}V = \frac{\partial }{\partial t} \left( \rho \, \mathrm{d}x \, \mathrm{d}y \, \mathrm{d}z \right) = \frac{\partial \rho}{\partial t} \, \mathrm{d}x \, \mathrm{d}y \, \mathrm{d}z
.\]
Hence we obtain the differential form of the mass conservation law:
\begin{equation}\label{eq:coneq1}
  \frac{\partial \rho u}{\partial x} + \frac{\partial \rho v}{\partial y} + \frac{\partial \rho w}{\partial z} + \frac{\partial \rho}{\partial t} = 0
\end{equation}
\Cref{eq:coneq1} is often called the \textit{continuity equation}.

Since $\nabla $ is given by:
\[ 
\nabla = \hat{i} \frac{\partial }{\partial x} + \hat{j} \frac{\partial }{\partial y} + \hat{k} \frac{\partial }{\partial z}
\]
then
\[ 
\frac{\partial \rho u}{\partial x} + \frac{\partial \rho v}{\partial y} + \frac{\partial \rho w}{\partial z} = \nabla \cdot \rho \textbf{V}
.\]
Therefore the conservation of mass from \Cref{eq:coneq1} can also be written as:
\begin{equation}\label{eq:coneqvec}
  \nabla \cdot \rho \textbf{V} + \frac{\partial \rho}{\partial t} = 0
\end{equation}
When \Cref{eq:coneq1} is written on vector form as in \Cref{eq:coneqvec} it may be applied in any coordinate system -- even a cylindrical or polar one.

Two special cases of the continuity equation are worth a mention.

\subsubsection{Incompressible Fluid}
For an incompressible fluid $\rho = \mathrm{constant}$ and is therefore neither a function of space coordinate nor time. Therefore the continuity equation for an incompressible fluid simplifies to:
\[ 
  \nabla \cdot  \textbf{V} = 0
.\]
Thus the velocity field, $\textbf{V}(x,y,z,t)$ must satisfy $\nabla \cdot \textbf{V} = 0$ for an incompressible fluid.


\subsubsection{Steady Flow}
For steady flow, all fluid properties, are, by definition, independent of time. This for steady flow the continuity equation simplifies to:
\[ 
\nabla \cdot \rho \textbf{V} = 0
.\]


\subsection{Stream Function for Two-Dimensional Incompressible Flow}
Streamlines have previously been defined as being lines tangent to the velocity vectors in a flow at an instant:
\[ 
\frac{\mathrm{d}y}{\mathrm{d}x} \bigg|_{\mathrm{streamline}} = \frac{v}{u}
.\]

We will now seek to develop a more formal definition of stream lines by introducing the \textit{stream function} $\psi$. This function allows us to represent two entities -- the velocity components $u(x,y,t)$ and $v(x,y,t)$ of a two-dimensional incompressible flow -- with a single function $\psi(x,y,t)$.

We start with a two-dimensional version of the continuity equation for incompressible flow:
\begin{equation}\label{eq:twodimconinc}
  \frac{\partial u}{\partial x} + \frac{\partial v}{\partial y} = 0
\end{equation}
We now define the stream function by:
\[ 
u \equiv \frac{\partial \psi}{\partial y} \qquad \mathrm{and} \qquad v \equiv - \frac{\partial \psi}{\partial x}
.\]
By doing this \Cref{eq:twodimconinc} is automatically satisfied for any $\psi(x,y,t)$. Using the ``naive'' definition of streamlines introduced at the beginning of this section we can obtain an equation that is only valid along a streamline:
\[ 
 u \, \mathrm{d}y  - v \, \mathrm{d}x  = 0
\]
or using the stream function
\begin{equation}\label{eq:strfunconeq2}
  \frac{\partial \psi}{\partial x} \, \mathrm{d}x + \frac{\partial \psi}{\partial y} \, \mathrm{d}y  = 0
\end{equation}

From a strictly mathematical point of view, at any instant in time $t$ the variation in a function $\psi(x,y,t)$ in space $(x,y)$ is given by;
\begin{equation}\label{eq:strfunmat}
  \mathrm{d}\psi = \frac{\partial \psi}{\partial x} \, \mathrm{d}x + \frac{\partial \psi}{\partial y} \, \mathrm{d}y 
\end{equation} 
Comparing \Cref{eq:strfunconeq2} and \Cref{eq:strfunmat} we see that along an instantaneous streamline $\mathrm{d}\psi = 0$, i.e. $\psi$ is constant along a streamline. Therefore we can specify individual streamlines by their stream function values $\psi = 0, 1, 2,\ldots $. Furthermore, these $\psi$ values can be used to obtain the volume flow rate between any two streamlines. Consider the streamlines shown on \Cref{fig:f5_3}. We can compute the volume flow rate between $\psi_1$ and $\psi_2$ by using line $AB$, $BC$, $DE$, or $EF$ (remember there is no flow across a streamline). 

\begin{figure} [ht]
  \centering
  \includegraphics[width=0.7\linewidth]{./figures/f5_3.png}
  \caption{Instantaneous streamlines in a two-dimensional flow}
  \label{fig:f5_3}
\end{figure}

Now let us compare the volume flow rate found using line $AB$ and line $BC$ -- these should be the same. 

For a unit depth, the flow rate across $AB$ is
\[ 
Q = \int_{y_1}^{y_2} u \, \mathrm{d}y = \int_{y_1}^{y_2} \frac{\partial \psi}{\partial y} \, \mathrm{d}y 
.\]
Along $AB$, $x = \mathrm{constant}$, thus
\[ 
\mathrm{d}\psi = \frac{\partial \psi}{\partial y} \, \mathrm{d}y 
\]
therefore:
\[ 
  Q = \int_{y_1}^{y_2} \frac{\partial \psi}{\partial y} \, \mathrm{d}y = \int_{\psi_1}^{\psi_2} \, \mathrm{d}\psi = \psi_2 - \psi_1
.\]

The same procedure is followed for $BC$, however here $y = \mathrm{constant}$, thus:
\begin{align*}
  Q &= \int_{x_1}^{x_2} v \, \mathrm{d}x  \\
  &= - \int_{x_1}^{x_2} \frac{\partial \psi}{\partial x} \, \mathrm{d}x  \\
  &= - \int_{\psi_2}^{\psi_1} \, \mathrm{d}\psi = \psi_2 - \psi_1
.\end{align*}
Hence the volume flow rate per unit depth between two streamlines is given by the difference between the stream function values.

For two-dimensional steady compressible flow in the $xy$ plane, the stream function $\psi$ can be defined such that:
\[ 
\rho u \equiv \frac{\partial \psi}{\partial y} \qquad \mathrm{and} \qquad \rho v \equiv - \frac{\partial \psi}{\partial x}
.\]
Then the difference between the constant values of $\psi$ defining two streamlines is the mass flow rate per unit depth between the two streamlines. 

If $\psi = 0$ is used to designate the streamline through the origin, the all other $\psi$ values for all other streamlines represents the flow between the origin and that streamline.
